\chapter{Shell Scripting and Command-Line Automation}\label{ch:shell}
\chapterguide{Chapter 2: files, command-line utilities, and exit status.}{use variables, quoting, redirection, pipelines, conditions, and loops to automate repeatable OS observations in Lab 3.}

\section{Why the shell belongs in an operating-systems course}
A shell such as Bash is an ordinary user-space program, but it is unusually useful for observing an operating system. It starts other programs, connects their input and output, reports whether they succeeded, and gives us a small programming language for repeating tasks.

This makes the shell a convenient bridge between theory and experiment. Instead of writing C code every time we want to inspect a process or count log messages, we can combine existing tools such as \texttt{ps}, \texttt{grep}, \texttt{wc}, \texttt{df}, and \texttt{uname}.

\begin{keyidea}
A good shell script is usually not doing everything itself. It coordinates small programs and lets the operating system connect them through processes, files, and standard streams.
\end{keyidea}

\section{What happens when you run a command}
When you enter an external command such as \texttt{ps}, the shell performs a simple sequence conceptually:
\begin{enumerate}
  \item interpret the command line and expand variables;
  \item prepare redirection or pipes if requested;
  \item start the requested program;
  \item wait for it if it is a foreground command;
  \item collect its exit status.
\end{enumerate}
You do not need to implement these steps in Bash, but recognising them helps connect shell work to processes, file descriptors, \texttt{fork()}, and \texttt{exec()}.

\section{Variables, arguments, and quoting}
A shell variable stores text. Script arguments are available as \texttt{\$1}, \texttt{\$2}, and so on. Quoting matters because the shell processes text before the final command runs.

\begin{lstlisting}[style=utpbash]
file=${1:-report.txt}
echo "Writing report to: $file"
df -h . > "$file"
\end{lstlisting}

The double quotes around \texttt{"\$file"} preserve a filename containing spaces as one argument.

\begin{pitfall}
Unquoted variables are a common source of scripting errors. If \texttt{file="OS Lab.txt"}, then \texttt{cat \$file} can be interpreted as two filenames. Prefer \texttt{cat "\$file"}.
\end{pitfall}

\section{Standard input, output, and error}
Most command-line programs begin with three standard streams:
\begin{itemize}
  \item \textbf{standard input} --- where normal input is read;
  \item \textbf{standard output} --- where normal results are written;
  \item \textbf{standard error} --- where diagnostics are written.
\end{itemize}
The shell can redirect these streams without changing the program.

\begin{utptable}
\renewcommand{\arraystretch}{1.18}
\begin{tabularx}{\linewidth}{@{}>{\bfseries\leavevmode\color{UTPNavy}}m{0.10\linewidth}
  >{\raggedright\arraybackslash}X
  >{\raggedright\arraybackslash}m{0.38\linewidth}@{}}
\toprule
\rowcolor{UTPPaleBlue!92}
\utpth{Syntax} & \utpth{Meaning} & \utpth{Example} \\
\midrule
\texttt{>} & replace a file with standard output & \texttt{ps > processes.txt} \\
\texttt{>>} & append standard output & \texttt{date >> history.txt} \\
\texttt{<} & take standard input from a file & \texttt{wc -l < events.log} \\
\texttt{2>} & redirect standard error & \texttt{cmd 2> errors.txt} \\
\texttt{|} & connect one program's output to another's input & \texttt{ps | grep bash} \\
\bottomrule
\end{tabularx}
\end{utptable}

\section{Pipelines are process composition}
Consider:
\begin{lstlisting}[style=utpbash]
ps -eo pid,ppid,stat,comm | grep bash
\end{lstlisting}
The first process produces text. The shell connects that output to a pipe. The second process reads from the pipe and keeps selected lines. The programs do not need to know each other's implementation.

\begin{intuition}
A pipeline is like an assembly line for bytes. Each program performs one small job, while the shell and operating system arrange the connections.
\end{intuition}

\section{Exit status is part of the interface}
Every command returns an integer status. By convention, 0 means success and a non-zero value means failure or a false result. The shell stores the most recent status in \texttt{\$?}.

\begin{lstlisting}[style=utpbash]
grep kernel notes.txt
echo $?
\end{lstlisting}
This is why commands can be composed inside \texttt{if}, \texttt{\&\&}, and \texttt{||} expressions.

\begin{workedexample}
A search tool can use status 0 for ``match found'', 1 for ``no match'', and 2 for ``could not read the file''. A script can then react to the condition without trying to interpret printed sentences.
\end{workedexample}

\section{Conditions and loops for repeatable observations}
Shell conditions and loops are most useful here when they automate repeated OS observations. For example, a script can count several log levels using one loop:
\begin{lstlisting}[style=utpbash]
for level in INFO WARN ERROR; do
    count=$(grep -c "^$level " events.log || true)
    echo "$level: $count"
done
\end{lstlisting}
The mathematics is simple; the important idea is repeatability. If the log changes, rerunning the same script applies the same method again.

\section{Background work and process observation}
Adding \texttt{\&} runs a command in the background:
\begin{lstlisting}[style=utpbash]
sleep 30 &
ps -o pid,ppid,state,comm
wait
\end{lstlisting}
This provides an easy way to create a temporary process and then inspect it. The \texttt{wait} command asks the shell to wait for its outstanding background job.

\section{A professional scripting workflow}
For a student lab script, a reliable workflow is:
\begin{enumerate}
  \item define the intended input and output;
  \item make one normal case work;
  \item quote variables and check required arguments;
  \item add one failure or boundary test;
  \item return a meaningful exit status;
  \item add a short usage message if another person will run the script.
\end{enumerate}

\begin{keyidea}
The goal of Bash in this course is not to memorise syntax. It is to make OS observations reproducible: the same commands, in the same order, with failures handled clearly.
\end{keyidea}

\begin{quickcheck}
\begin{enumerate}
  \item Why should a filename variable usually be quoted?
  \item What OS idea is visible in a pipeline such as \texttt{ps | grep bash}?
  \item Why is exit status more useful to another program than a human-readable sentence alone?
  \item What does appending \texttt{\&} change about how the shell treats a command?
\end{enumerate}
\end{quickcheck}
